% !TeX root = ../cosmology_summary.tex

%% --------------------------------------------------------
\chapter{Гідродинаміка фотон-баріонної плазми\\ та акустичні осциляції}
\label{sec:bao}
%% --------------------------------------------------------

Після розігріву Всесвіт входить у радіаційно-доміновану епоху.
Температура настільки висока, що баріонна матерія повністю
іонізована: протони та електрони утворюють плазму, в якій
фотони інтенсивно розсіюються на вільних електронах
(томсонівське розсіяння). Середня довжина вільного пробігу
фотона набагато менша за космологічний горизонт, тому
випромінювання і баріони поводяться як єдина
\emph{фотон-баріонна плазма} \cite{mukhanov_2005}.

Первинні неоднорідності кривини $\mathcal{R}_k$, успадковані
від інфляції, стають початковими умовами для цієї плазми.
Подальша динаміка визначається конкуренцією двох сил:
гравітаційне притягання намагається стиснути речовину у
потенціальні ями $\Phi$, а внутрішній тиск фотонного газу
чинить опір. Результат~--- акустичні коливання.

%% --------------------------------------------------------
\section{Рівняння акустичного осцилятора}
%% --------------------------------------------------------

Щоб описати ці коливання кількісно, потрібно знати
швидкість звуку в плазмі.
Швидкість звуку визначається як $c_s^2 = \partial p/\partial\rho$.
Оскільки тиск цілком фотонний ($p_\gamma = \bar{\rho}_\gamma/3$,
а баріонний тиск нехтовно малий~--- нерелятивістський протонний
газ при $T \sim 10^4$~К має $c_s^{(b)} \sim \sqrt{T/m_p}
\sim 10^{-5}\,c$), то $c_s = 1/\sqrt{3} \approx 0.577\,c$
у чистій фотонній плазмі.

В релятивістській гідродинаміці маса спокою --- не єдиний
внесок в інерцію: тиск теж чинить опір прискоренню, і
роль \enquote{ефективної маси} відіграє ентальпія $\rho + p$.
Для фотонів $\rho_\gamma + p_\gamma = \tfrac{4}{3}\rho_\gamma$~---
саме це число визначає, наскільки «важкою» є фотонна компонента
як партнер баріонів у спільних коливаннях. Зручно ввести
безрозмірний параметр, що показує відносну вагу баріонної
інерції порівняно з фотонною (baryon loading parameter):
\begin{equation*}
  R(\eta) \equiv
    \frac{\bar{\rho}_b}{\rho_\gamma + p_\gamma}
    = \frac{3\bar{\rho}_b(\eta)}{4\bar{\rho}_\gamma(\eta)}
    \propto a(\eta).
\end{equation*}
$R$ зростає з часом, бо $\bar{\rho}_b \propto a^{-3}$
спадає повільніше за $\bar{\rho}_\gamma \propto a^{-4}$.
Фізична аналогія: додаємо масу до маятника без зміни
пружини~--- частота коливань падає.
Підставляючи $R$ у вираз для швидкості звуку:
\begin{equation}\label{eq:sound_speed}
  c_s(\eta) = \frac{1}{\sqrt{3\bigl(1 + R(\eta)\bigr)}}.
\end{equation}
При $R \to 0$ отримуємо $c_s \to 1/\sqrt{3}$~--- чиста
фотонна плазма; при $R \gg 1$ (пізня епоха, багато баріонів)
$c_s \to 0$.

Акустичні коливання фізично можливі лише на субхабблових
масштабах ($k\eta \gg 1$): тільки тоді різні точки плазми
перебувають у причинному контакті і можуть коливатися
узгоджено.

Виведення рівняння осцилятора починається з двох рівнянь
гідродинаміки у Фур'є-просторі для збурення фотонної густини
$\delta_\gamma \equiv \delta\rho_\gamma/\bar{\rho}_\gamma$
та дивергенції поля швидкостей плазми $\theta \equiv \nabla\cdot\mathbf{v}$
\cite{dodelson_2020}:

\textbf{Рівняння неперервності} (збереження фотонів):
\begin{equation*}
  \delta_\gamma' = -\frac{4}{3}\,\theta + 4\Phi',
\end{equation*}
де $4/3 = 3(1+w)$ при $w=1/3$~--- це стандартний коефіцієнт
рівняння збереження $\dot\rho + 3H(1+w)\rho = 0$ для відносного
збурення фотонної густини,
а $\Phi'$ — через зміну об'єму в збуреній метриці.

\textbf{Рівняння Ейлера} (баланс тиску і гравітації) для
об'єднаної фотон-баріонної рідини з ефективною інерцією $(1+R)$:
\begin{equation*}
  \theta' + \mathcal{H}\frac{R}{1+R}\,\theta
  = \frac{k^2}{4(1+R)}\,\delta_\gamma + k^2\Psi.
\end{equation*}
Член $\mathcal{H}R/(1+R)$ — це хабблівське тертя, ефективне
лише для баріонної компоненти ($R$); фотони не відчувають тертя.

Диференціюємо рівняння неперервності за $\eta$ і підставляємо
$\theta'$ з рівняння Ейлера. У конформно-ньютонівському
калібруванні $\Phi = \Psi$ і на субхабблових масштабах
членами з $\mathcal{H}$ нехтуємо. Після скорочень отримуємо
рівняння вимушеного осцилятора \cite{dodelson_2020}:
\begin{equation}\label{eq:acoustic_oscillator}
  \frac{d}{d\eta}\bigl[(1+R)\,\delta_\gamma'\bigr]
  + \frac{k^2}{3}\,\delta_\gamma
  = -\frac{k^2}{3}(1+R)\,\Phi.
\end{equation}
Ліва частина~--- осцилятор зі змінною масою $(1+R)$
і частотою $\omega = kc_s$; коефіцієнт $1/3$ є саме
$c_s^2 = 1/(3(1+R))$ з~\eqref{eq:sound_speed}. Права
частина~--- гравітаційне вимушення потенціалом $\Phi$,
що зміщує положення рівноваги. Кожна мода $k$ коливається
незалежно, а оскільки всі вони стартують з однакової
початкової фази $\mathcal{R}_k$~--- на відміну від
некогерентних джерел типу топологічних дефектів~--- їхні
осциляції є \emph{когерентними}: моди з однаковим $k$
перебувають в одній фазі незалежно від просторового
положення. Саме когерентність породжує чіткі акустичні
піки на спектрі CMB.

%% --------------------------------------------------------
\section{Рекомбінація, звуковий горизонт і BAO}
%% --------------------------------------------------------

Акустичні коливання тривають до \emph{рекомбінації}
($z_\mathrm{rec} \approx 1100$, $T \approx 3000$~К):
протони захоплюють електрони, іонізація припиняється,
томсонівське розсіяння вимикається і фотони починають
вільно поширюватися. Плазма «заморожується» в тому стані,
в якому її застала рекомбінація.

Максимальна відстань, яку встигла пройти звукова хвиля
від початку гарячої епохи до рекомбінації,~---
\emph{радіус звукового горизонту}:
\begin{equation}\label{eq:sound_horizon}
  r_s(\eta_\mathrm{rec})
    = \int_0^{\eta_\mathrm{rec}} c_s(\eta)\,d\eta
    \approx 147\,\mathrm{Mpc}.
\end{equation}
Це єдиний лінійний масштаб у задачі. Моди, що в момент
рекомбінації перебували у точці максимального стиснення
або розширення, проявляються у температурному спектрі
CMB як серія \emph{акустичних піків}~--- докладніше
у §\,\ref{sec:sachs}.

Але акустичні хвилі залишають слід не лише у фотонах.
Після рекомбінації фотони пішли геть, а баріонна речовина
«замерзла» у вигляді сферичної оболонки радіусом $r_s$
навколо кожної початкової неоднорідності темної матерії
(рис.~\ref{fig:bao_profile}). Це~--- \textbf{баріонні
акустичні осциляції (BAO)} \cite{eisenstein_2005}.

\begin{figure}[ht!]
\centering
\begin{tikzpicture}
  \begin{axis}[
    width=\linewidth,
    height=0.75\linewidth,
    xmin=0, xmax=250,
    ymin=-0.05, ymax=1.1,
    xlabel={Відстань від первинного збурення, $r$ [Мпк]},
    ylabel={Контраст густини баріонів,
            $\Delta\rho_b/\bar{\rho}_b$},
    grid=both,
    grid style={line width=.1pt, color=gray!20},
    major grid style={line width=.2pt, color=gray!40},
    tick label style={font=\footnotesize},
    label style={font=\small},
    legend style={at={(0.95,0.95)}, anchor=north east,
                  font=\footnotesize, cells={anchor=west}},
    axis lines=left,
    clip=false
  ]
    \addplot[
      domain=0:240, samples=200,
      color=accent, line width=1.5pt
    ] {exp(-0.015*x) + 0.35*exp(-((x-147)/12)^2)};
    \addlegendentry{Профіль густини речовини}

    \draw[dashed, color=theoryred, line width=1pt]
      (axis cs:147,-0.05) -- (axis cs:147,0.9);
    \node[anchor=south, color=theoryred, font=\footnotesize]
      at (axis cs:147,0.9) {$r_s \approx 147\text{ Мпк}$};

    \node[anchor=west, color=gray, font=\scriptsize]
      at (axis cs:165,0.2) {Затухання Сілка};
    \draw[->, color=gray, >=stealth]
      (axis cs:162,0.18) -- (axis cs:155,0.1);
  \end{axis}
\end{tikzpicture}
\caption{Схематичний профіль поширення первинного
  адіабатичного збурення у баріонній компоненті плазми
  (модельна крива). Надлишок речовини на радіусі
  звукового горизонту $r_s$ відповідає сучасному
  масштабу BAO.}
\label{fig:bao_profile}
\end{figure}

У сучасну епоху масштаб BAO проявляється як фіксований
надлишок у просторовій кореляційній функції розподілу
галактик на відстані $\sim 147$~Мпк. Оскільки $r_s$
жорстко визначений лінійною фізикою раннього
Всесвіту~\eqref{eq:sound_horizon}, BAO~--- ідеальна
«стандартна лінійка»: вимірюючи цей масштаб на різних
червоних зміщеннях, можна реконструювати $H(z)$ і
параметри темної енергії незалежно від CMB.

Той самий звуковий горизонт $r_s$, що визначає масштаб
BAO у розподілі галактик, відповідає і за положення
акустичних піків у температурному спектрі CMB. Як саме
ці піки формуються~--- і які інші ефекти на них
накладаються~--- розглядається у наступному розділі.

